\documentclass[11pt]{article}
\usepackage[margin=1in]{geometry}
\usepackage{amsmath,amssymb,amsthm,booktabs,array}
\usepackage[hidelinks]{hyperref}
\newtheorem{theorem}{Theorem}
\newtheorem{lemma}{Lemma}
\newtheorem{proposition}{Proposition}
\newcommand{\B}{B_6}
\newcommand{\Mod}{\operatorname{Mod}}
\newcommand{\Aut}{\operatorname{Aut}}
\newcommand{\normal}[1]{\langle\!\langle #1\rangle\!\rangle}
\title{Fixed-generator positive factorizations in Wajnryb's Artin \(A_5\) quotient}
\author{Alec Kriebel\\
\small Independent researcher\\
\small ORCID: \href{https://orcid.org/0009-0001-9320-500X}{0009-0001-9320-500X}}
\date{October 3, 2026}
\begin{document}
\maketitle
\begin{abstract}
We classify the positive words in the five standard generators of
\(B_6/\normal{ch^{-1}}\) representing \(c^2\), where
\(c=(a_1a_2a_3a_4)^5\) and \(h=a_5a_4a_3a_2a_1^2a_2a_3a_4a_5\).
Their lengths are \(20,30,40\), with respectively one, one, and two
strict Hurwitz classes. Simultaneous conjugation identifies the two
length-\(40\) classes, giving the three representatives proposed in
Wajnryb's question. The proof combines classical geometric and braid
results with pure-braid linking constraints and an exhaustive,
exact \(810\)-word certificate for the length-\(20\) case.
The length spectrum and the twice-per-pair crossing theorem used at
length \(30\) have prior sources; a reproducible six-strand certificate
for the latter is supplied as additional evidence.
\end{abstract}

\section{The question and equivalence conventions}

Let \(a_1,\ldots,a_5\) be the standard generators of the braid group
\(\B\), with \(a_i a_{i+1}a_i=a_{i+1}a_i a_{i+1}\) and
\(a_i a_j=a_j a_i\) for \(|i-j|>1\). Put
\[
 c=(a_1a_2a_3a_4)^5,\qquad
 h=a_5a_4a_3a_2a_1^2a_2a_3a_4a_5,\qquad
 z=(a_1a_2a_3a_4a_5)^6,\qquad
 G=\B/\normal{ch^{-1}}.
\]
The classical full-twist decomposition \(z=ch=hc\) implies
\begin{equation}\label{central}
 c^2=h^2=z \quad\text{in }G;
 \qquad z\text{ is central.}
\end{equation}
In \cite[Section 3, p.~126 of the linked online version]{W},
Wajnryb asks whether every positive word representing \(c^2\) in \(G\)
has length \(20,30,\) or \(40\), and is Hurwitz equivalent to one of
the three displayed models \(h^2,z,c^2\).

Here a \emph{positive word} means a literal word in these five fixed
generators. Its factorization is the ordered tuple of its letters
in \(G\). A strict Hurwitz move replaces adjacent factors
\((x,y)\) by \((xyx^{-1},x)\), or applies its inverse.
We separately consider the equivalence generated by these moves
and simultaneous conjugation of all factors by an element of \(G\).
Auroux introduces these operations separately in
\cite[Section 1, p.~131 of the linked online version]{A}.
Intermediate Hurwitz tuples may contain conjugates of generators;
the inputs quantified here remain fixed-generator words.

\begin{theorem}\label{main}
The positive words in \(a_1,\ldots,a_5\) representing \(c^2\) in \(G\)
have precisely the following strict Hurwitz classes:
\[
\begin{array}{ccl}
\text{length}&\text{number of classes}&\text{representatives}\\
\hline
20&1&h^2\\
30&1&(a_1a_2a_3a_4a_5)^6\\
40&2&(a_1a_2a_3a_4)^{10},\quad(a_2a_3a_4a_5)^{10}.
\end{array}
\]
Every length-\(20\) word in this set is literally one of the ten
distinct cyclic rotations of \(h^2\). With simultaneous conjugation,
there are exactly three classes, one at each length.
\end{theorem}

Thus the proposed three-model classification holds under the
equivalence including simultaneous conjugation. Under strict moves
alone the length-\(40\) model splits into two classes.
This theorem resolves the fixed-generator quotient question under
both conventions. It does not classify arbitrary products of
conjugates of generators or arbitrary nonseparating Dehn twists.

\paragraph{Prior results and contribution.}
The exact quotient is a classical genus-two, one-boundary mapping
class presentation of Matsumoto \cite[Theorem 1.3]{M}:
the relation \(z=c^2\) is equivalent to \(c=h\).
The numerical spectrum \(20,30,40\) follows already from
Sato's geography theorem \cite[Theorem 5.1, Table 1, Remark 5.1(1)]{S}
applied to the associated relatively minimal genus-two fibration
with a section of square \(-1\) and no reducible singular fibers.
We use the later maximum-length theorem \cite[Theorem A]{BMVHM}
and an elementary linking argument to obtain that spectrum here.
Moreover, Elrifai--Morton \cite[p.~496]{EM} report Casson's general
result that a positive braid with exactly two crossings per labelled
pair is the full twist. They do not supply its proof or a separate
Casson reference. Our six-strand computation verifies this known
assertion in the required rank; it is not a new general theorem.
The point of the note is the complete fixed-input orbit
classification, especially the length-\(20\) reduction and the
strict versus simultaneous-conjugation distinction.
A literature audit located no publication of this exact complete
classification; that search does not establish absence of prior work.

\section{Length and linking restrictions}

\subsection{Geometric bound}
Take the standard five-curve chain on a genus-two surface with two
boundary components and cap one boundary. The five twists define
a homomorphism \(\B\to\Mod(\Sigma_{2,1})\). The odd and even chain
relations send \(z\) and \(c^2\) to the twist \(t_\delta\) about the
remaining boundary. Consequently \(z=ch\) forces the images of
\(c\) and \(h\) to agree, so the homomorphism factors through \(G\).
All five curves are nonseparating. The genus-two, one-boundary case
of \cite[Theorem A]{BMVHM} says that the maximum number of
nonseparating positive twists in a factorization of \(t_\delta\)
is \(40\). Hence every word under consideration has length \(m\leq40\).
Only this homomorphism is needed for the bound, not an injectivity
assertion for \(G\).

The abelianization of \(\B\) is infinite cyclic, all \(a_i\) having
the same class. The added relation has exponent sums \(20\) and
\(10\), so \(G_{\rm ab}\cong\mathbb Z/10\mathbb Z\).
The map sending each \(a_i\) to \(1\) modulo \(10\) proves that
there are no additional abelian relations.
Thus \(m\in\{0,10,20,30,40\}\).

\subsection{Normal-closure constraint}
Let \(P_6\) be the pure braid group. Both \(c\) and \(h\) are pure;
the permutation map to \(S_6\) therefore descends to \(G\).
A word representing \(c^2\) is pure as a braid in \(\B\).
For a pure braid let \(\ell_{ij}\) be half the signed crossing
count of labelled strands \(i,j\). These are integral homomorphisms
\(P_6\to\mathbb Z\), and braid conjugation permutes their coordinates.
For a positive pure word all \(\ell_{ij}\geq0\) and
\(m=2\sum_{i<j}\ell_{ij}\).

The linking vector \(\beta\) of \(c^2\) equals \(2\) on pairs among
\(\{1,\ldots,5\}\) and \(0\) on pairs incident to \(6\).
The linking vectors of conjugates of \(ch^{-1}\) are
\((r_k)_{ij}=1-2\,\mathbf1_{k\in\{i,j\}}\), \(1\leq k\leq6\).
Expressing \(w(c^2)^{-1}\) as a product of conjugates of the relator
and its inverse gives integers \(t_1,\ldots,t_6\), of either sign,
with
\begin{equation}\label{link}
 \ell_{ij}(w)=\beta_{ij}+T-2(t_i+t_j)\geq0,\qquad
 T=\sum_{i=1}^6t_i=m/10-4.
\end{equation}
This is a necessary constraint only; no sufficiency of linking
data is assumed.

\begin{lemma}\label{patterns}
Lengths \(0\) and \(10\) are impossible. At lengths \(20,30,40\),
the linking vectors are respectively twice a star, the all-ones
vector, and twice a complete graph on five vertices with one
isolated vertex.
\end{lemma}
\begin{proof}
Write \(s=\sum_{i=1}^5t_i\), so \(t_6=T-s\).
The ten internal and five remaining inequalities in \eqref{link} are
\[
 t_i+t_j\leq\lfloor(T+2)/2\rfloor\quad(i,j\leq5),\qquad
 t_i+t_6\leq\lfloor T/2\rfloor\quad(i\leq5).
\]
For \(T=-4,-3\) the internal sums yield \(s\leq-3\);
the remaining sums yield
\(s\geq\lceil(5T+10)/4\rceil=-2,-1\), respectively.
This excludes \(m=0,10\).

For \(T=-2\), summing gives \(-1\leq s\leq0\).
If \(s=-1\), then \(t_6=-1\), all \(t_i\leq0\) for \(i\leq5\),
and exactly one is \(-1\). If \(s=0\), then \(t_6=-2\).
Any positive maximum \(M\) among the first five entries would make
the other four at most \(-M\), giving \(s\leq-3M<0\).
Thus the first five entries are all zero.
Substitution gives precisely the six possible doubled stars.

For \(T=-1\), the summed inequalities force \(s=0,t_6=-1\);
the remaining inequalities make the first five entries at most zero,
so all are zero. Hence every \(\ell_{ij}=1\).

For \(T=0\), the remaining summed inequalities give \(t_6\leq0\),
so \(s\geq0\). If the maximum first entry \(M\geq1\), the other four
are at most \(1-M\), giving \(s\leq4-3M\leq1\);
if \(M\leq0\), then \(s\leq0\).
When \(s=0\), all entries vanish. When \(s=1\), \(t_6=-1\),
each first entry is at most \(1\), at most one can be \(1\),
and their sum forces one \(1\) and four zeros.
These give exactly the six doubled five-vertex cliques.
\end{proof}
The models in \eqref{central} realize all three permitted lengths.
This proves the length assertion of Theorem~\ref{main}.

\section{Length \texorpdfstring{\(40\)}{40}: two strict classes}
An isolated strand in the linking vector of a positive word has
no crossings. It stays at its initial position \(k\).
If \(2\leq k\leq5\), strands initially to its left and right cannot
cross one another, contradicting the prescribed clique.
Thus \(k=1\) or \(6\).

For \(k=6\), only \(a_1,\ldots,a_4\) occur.
The four-chain map \(B_5\to\Mod(\Sigma_{2,1})\) is injective by
the \(A_4\) case of Perron--Vannier, as stated in
\cite[Theorem 2.1]{W}. Its chain neighborhood is the genus-two,
one-boundary surface, up to a boundary collar.
The geometric images of \(w\) and \(c^2\) agree, so they are equal
already in \(B_5\).
The positive braid monoid embeds in \(B_5\) \cite[Section 4]{GM};
hence their words are connected by positive braid and commutation
relations.
A commutation is a Hurwitz swap. For \(xyx=yxy\), two inverse
Hurwitz moves give
\[
 (x,y,x)\longmapsto(y,y^{-1}xy,x)
 \longmapsto(y,x,x^{-1}y^{-1}xyx)=(y,x,y).
\]
Thus every such tuple is strictly equivalent to \(c^2\).
The other endpoint gives the class of
\(w_0=(a_2a_3a_4a_5)^{10}\).

With \(d=a_1a_2a_3a_4a_5\), braid relations give
\(da_i d^{-1}=a_{i+1}\) for \(i\leq4\), so
\(dc^2d^{-1}=w_0=z\) in \(G\).
The permutation images of the factor-generated subgroups for the
two models are the \(S_5\) fixing \(6\) and the \(S_5\) fixing \(1\).
A strict Hurwitz move preserves the subgroup generated by all
factors, so these models are strictly inequivalent.
Simultaneous conjugation by \(d\) identifies them.

\section{Length \texorpdfstring{\(20\)}{20}: a complete finite certificate}
Lemma~\ref{patterns} forces one labelled strand to cross each
of the other five exactly four times; none of the other strands
cross each other. Their relative order is constant.
The distinguished strand makes a nearest-neighbor walk on six
positions, returns to its start, and uses every edge four times.
Crossing edge \(i\) records exactly the letter \(a_i\).
The fixed order of the other strands makes edge crossings
correspond to crossings with individual labels.
Conversely every such walk gives this pure-braid linking pattern.

Enumerate walks recursively, choosing either adjacent position
whenever available and its edge has been used fewer than four
times. At depth \(20\), retain a walk exactly when it returns to
its starting position. The five edge caps force all edge counts
to be four at that depth. This explores every candidate without
any other pruning.
The counts for starts \(1,\ldots,6\) are
\[
 81,\ 162,\ 162,\ 162,\ 162,\ 81,
 \qquad \text{total }810.
\]

We now give a presentation-compatible rejection test.
Let \(F_4=\langle x_1,x_2,x_3,x_4\rangle\).
In the following table, a signed integer string is a free word:
\(12\) denotes \(x_1x_2\), and \(\bar j\) denotes \(x_j^{-1}\).
Use right actions, applying the first word and then the second.
\[
\begin{array}{c|cccc}
 &x_1&x_2&x_3&x_4\\\hline
A_1&1&12&13&14\\
A_2&1\bar2 1&1&3&4\\
A_3&1&2\bar3 2&2&4\\
A_4&1&2&3\bar4 3&3\\
A_5&1&2&3&3\bar2 14\\\hline
A_1^{-1}&1&\bar1 2&\bar1 3&\bar1 4\\
A_2^{-1}&2&2\bar1 2&3&4\\
A_3^{-1}&1&3&3\bar2 3&4\\
A_4^{-1}&1&2&4&4\bar3 4\\
A_5^{-1}&1&2&3&\bar1 2\bar3 4
\end{array}
\]
Substitution and free reduction verify both inverse compositions,
the four adjacent braid relations, six distant commutations,
and \(A_c=A_h\). These finite word equalities prove that the
tables define a right action of the presented group \(G\).
No faithfulness is asserted or required.

Apply this action to each of the \(810\) words and compare all
four reduced images with those of \(c^2\).
Exactly ten words pass; the passing set is literally the set of
distinct cyclic rotations of \(h^2\).
The numbers passing for the six starts are \(1,2,2,2,2,1\).
These conclusions are exhaustive computation statements:
the supplement contains the complete automorphism checks,
all \(810\) word/image records, the survivor list, and independent
implementations. All comparisons use full reduced words, not hashes.

Rejection is sound because the action is defined on \(G\).
Acceptance has a separate group certificate: \(h^2=z\) is central,
so every cyclic rotation of \(h^2\) has product \(z\).
For a tuple with central product, successive inverse Hurwitz moves
moving its first factor to the end yield the ordinary cyclic rotation:
the final conjugated factor is \(z^{-1}g_1z=g_1\).
Thus the ten words form exactly one strict Hurwitz class.

\section{Length \texorpdfstring{\(30\)}{30} and the reproducible crossing certificate}
The linking vector is all ones, so every labelled pair crosses
exactly twice. Casson's result, reported in \cite[p.~496]{EM},
would give \(w=z\) in \(\B\).
For a checkable proof in the needed six-strand case, we instead
give a complete finite certificate of that equality.

Index the fifteen unordered label pairs lexicographically.
A state consists of crossing counts \(q_{ij}\in\{0,1,2\}\)
and the order \(p\) of the six strand labels.
Encode the counts by \(q=\sum q_{ij}3^{\operatorname{index}(i,j)}\).
Initially \(q=0\) and \(p\) is the identity.
Each adjacent positive generator swaps the adjacent labels and
increments their pair count if it is less than two.
Counts increase by one, so the graph is a finite directed acyclic
graph of depth at most \(30\).

The parity of each pair count determines whether that pair's
relative order has reversed. Thus a reachable count vector
determines a unique permutation, and merging by count code
does not lose continuations. Breadth-first search exhausts
all legal edges and yields \(234{,}368\) reachable states \(R\)
and \(711{,}342\) outgoing edges.
The full vector \(f\), all counts two, is reachable with identity
permutation. Complete candidate words are exactly paths from \(0\)
to \(f\).

A reachable vector \(q\) can continue to \(f\) exactly when
\(f-q\) is reachable. Indeed, reverse the sequence of swaps in
a completing suffix, keeping generators positive: it starts at
the identity and has labelled counts \(f-q\).
Conversely a reaching word for \(f-q\) ends at the permutation
of \(q\), since the two vectors have the same parities.
Reversing that sequence gives a legal completion from \(q\);
nonnegative added counts never exceed the final counts.
This is reversal of a positive word, not a braid inverse.
Hence the coaccessible subgraph
\[
 S=\{q\in R:f-q\in R\}
\]
has \(90{,}921\) states and \(261{,}810\) edges.
Independent backward exploration from \(f\) gives the same set.

Use the faithful Artin action on \(F_6\), in the right-action convention
\[
 T_i(x_i)=x_i x_{i+1}x_i^{-1},\qquad
 T_i(x_{i+1})=x_i,\qquad T_i(x_j)=x_j\quad(j\ne i,i+1).
\]
Faithfulness is classical \cite[Section 1.6]{GM}; the inverse-generator
and composition conventions change it by bijective braid
automorphisms or anti-automorphisms, so still detect equality.
Choose a reaching parent path to each state. A parent of a state
in \(S\) lies in \(S\), since it inherits a completion.
Let \(U(q)\) be the exact six reduced free-word images along this
path. The certificate verifies, on \emph{every} edge
\(q\xrightarrow{i}q'\) in \(S\),
\begin{equation}\label{edge}
 T_i\bigl(U(q)\bigr)=U(q')
\end{equation}
by full word equality, and verifies that \(U(f)\) equals the
action of the displayed word \(z\).
Induction along an arbitrary complete path now gives its action
as \(U(f)\). Faithfulness proves \(w=z\) in \(\B\).
Positive monoid embedding and the Hurwitz realization of braid
relations then put all length-\(30\) words in one strict class.

Counting states or equating hashes would not prove \eqref{edge}.
Python and C++ implementations perform all \(261{,}810\) literal
edge comparisons. The distributed controls
\texttt{independent\_20\_30.py} and
\texttt{independent\_ranks1\_6.py} each reconstruct the backward
graph and reproduce them. Their complete sorted \(90{,}921\)-record action
streams agree byte for byte.
For integrity, the record stream has \(15{,}258{,}431\) bytes and
SHA-256
\begin{center}\small\ttfamily
af2b8ec569d613e4f3d8ba3b72d18e85e8c6\\
125f4265057b7f4c485d544d159bf
\end{center}
Equal pair counts alone do not determine a braid: the positive
prefixes \(a_1^2a_2^2\) and \(a_2^2a_1^2\) have equal pair counts
and permutations but different free-group actions.
They lie outside \(S\), and the independent negative controls
verify this distinction.

\section{Verification, scope, and disclosure}
The supplement distributes the original four-turn proofs and
checkers, exact receipts, complete compressed action records,
independent finite controls, and a scope-specific priority report.
The length-\(20\) test requires only Python's standard library;
the second large certificate additionally uses a C++17 compiler.
No external group solver or numerical approximation is needed.
Source hashes and package hashes are integrity checks;
the completeness arguments and exact comparisons above remain
the mathematical content of the certificates.
The geometric length theorem, the chain embedding, faithful
Artin action, and positive monoid embedding are credited
classical inputs rather than claims proved by these programs.

Different lengths cannot be Hurwitz equivalent.
At length \(40\) the generated-subgroup invariant proves that
the two strict classes are distinct; the explicit conjugation
merges them. The other two lengths each have one strict class.
This completes the proof of Theorem~\ref{main}.

\paragraph{Use of AI and review status.}
AI tools were used extensively in developing the proofs,
writing and running verification code, searching the literature,
drafting this note, and conducting adversarial reviews.
Independent AI review paths are recorded in the accompanying
research audit. This manuscript is an unrefereed preprint;
it has not received external human peer review.
The stated result concerns only the fixed standard-generator
input universe and both expressly defined equivalence conventions.

\begin{thebibliography}{9}
\bibitem{A}
D. Auroux, \emph{Mapping class group factorizations and symplectic
4-manifolds: some open problems}, in B. Farb (ed.),
Problems on Mapping Class Groups and Related Topics,
Proc. Sympos. Pure Math. 74, AMS, 2006, Chapter 9.
\href{https://www.math.uchicago.edu/~farb/papers/mcgbook.pdf}{Online version}.

\bibitem{BMVHM}
R. I. Baykur, N. Monden, and J. Van Horn-Morris,
\emph{Positive factorizations of mapping classes},
Algebr. Geom. Topol. 17 (2017), 1527--1555.
\href{https://doi.org/10.2140/agt.2017.17.1527}{doi:10.2140/agt.2017.17.1527}.

\bibitem{EM}
E. A. Elrifai and H. R. Morton,
\emph{Algorithms for positive braids},
Quart. J. Math. Oxford (2) 45 (1994), 479--497.
\href{https://doi.org/10.1093/qmath/45.4.479}{doi:10.1093/qmath/45.4.479}.

\bibitem{GM}
J. Gonz\'alez-Meneses,
\emph{Basic results on braid groups},
Ann. Math. Blaise Pascal 18 (2011), 15--59.
\href{https://www.numdam.org/article/AMBP_2011__18_1_15_0.pdf}{Official text}.

\bibitem{M}
M. Matsumoto, \emph{A presentation of mapping class groups in terms
of Artin groups and geometric monodromy of singularities},
Math. Ann. 316 (2000), 401--418.
\href{https://doi.org/10.1007/s002080050336}{doi:10.1007/s002080050336}.
Original manuscript MPIM1997-74,
\href{https://archive.mpim-bonn.mpg.de/id/eprint/3423/}{institutional archive}.

\bibitem{S}
Y. Sato, \emph{2-Spheres of Square \(-1\) and the Geography of
Genus-2 Lefschetz Fibrations},
J. Math. Sci. Univ. Tokyo 15 (2008), 461--491.
\href{https://www.ms.u-tokyo.ac.jp/journal/pdf/jms150403.pdf}{Official text}.

\bibitem{W}
B. Wajnryb, \emph{Relations in the mapping class group},
in B. Farb (ed.), Problems on Mapping Class Groups and Related Topics,
Proc. Sympos. Pure Math. 74, AMS, 2006, Chapter 8.
\href{https://www.math.uchicago.edu/~farb/papers/mcgbook.pdf}{Online version}.
\end{thebibliography}
\end{document}
